\chapter{总复习}

\Needspace{12\baselineskip}

\section{总复习（一）}

\Needspace{8\baselineskip}

\pdfbookmark[2]{原书第81页}{orig-81}

\subsection*{原书第81页}

% p081_s01
\BookQuestion{一、1}
阳光文体超市正在进行促销活动，12根跳绳需要300元，每根跳绳是（\hspace{1.8em}）元，张老师准备带2600元去购买90根跳绳，能找回（\hspace{1.8em}）元。\par
\begin{BookAnswer}
依次填\(25\)、\(350\)。每根\(300\div12=25\)元；买\(90\)根用\(25\times90=2250\)元，找回\(2600-2250=350\)元。
\end{BookAnswer}

% p081_s02
\BookQuestion{一、2}
要使\(16\times(\square\div4)=512\)，\allowbreak{}\(\square\)里应该填（\hspace{1.8em}）。\par
\begin{BookAnswer}
填\(128\)。把16拆成\(4\times4\)，与括号内除以4抵消一个4；于是整个原式相当于所填数乘4。\[\begin{aligned}&512\div(16\div4)\\={}&512\div4\\={}&128\end{aligned}\]检验：\(16\times(128\div4)=16\times32=512\)。
\end{BookAnswer}

% p081_s03
\BookQuestion{一、3}
在精密电子测量中，一款芯片的长度由3个十、6个十分之一和9个千分之一组成，这个数写作（\hspace{1.8em}），读作（\hspace{1.8em}），保留两位小数是（\hspace{1.8em}）。\par
\begin{BookAnswer}
依次填\(30.609\)、三十点六零九、\(30.61\)。\(3\)个十是\(30\)，\allowbreak{}\(6\)个十分之一是\(0.6\)，\allowbreak{}\(9\)个千分之一是\(0.009\)，相加得\(30.609\)。百分位应写\(0\)占位，读数时这个\(0\)也要读。保留两位小数看千分位\(9\)，向百分位进\(1\)，故\(30.609\approx30.61\)。
\end{BookAnswer}

% p081_s04
\BookQuestion{一、4}
据统计，我国超过60岁的老人约331000000人，把横线上的数改写成用“亿”作单位的数是（\hspace{1.8em}）亿，省略“亿”位后面的尾数约是（\hspace{1.8em}）亿。\par
\begin{BookAnswer}
依次填\(3.31\)、\(3\)。改写成以亿为单位，在亿位右侧点上小数点，数的大小不变：\[331000000=3.31\text{亿}\]省略亿位后面的尾数，要看千万位的\(3<5\)，舍去，得到\[331000000\approx3\text{亿}\]前者是准确改写，用等号；后者是近似数，用约等号。
\end{BookAnswer}

% p081_s05
\BookQuestion{一、5}
小明由于粗心大意，把\((\triangle-25)\times12\)错算成了\(\triangle\times12-25\)，正确结果与他算的结果相差（\hspace{1.8em}）。\par
\begin{BookAnswer}
填\(275\)。正确算式减的是\(25\times12=300\)，错算只减\(25\)，所以错得数比正确得数多\(300-25=275\)。也可算\(25\times(12-1)=275\)。
\end{BookAnswer}

% p081_s06
\BookQuestion{一、6}
学校运动会上，小兵代表班级参加1500米赛跑，他平均每分钟跑240米，跑6分钟后，离终点还有（\hspace{1.8em}）米。\par
\begin{BookAnswer}
填\(60\)。六分钟跑\(240\times6=1440\)米，离终点\(1500-1440=60\)米。
\end{BookAnswer}

% p081_s07
\BookQuestion{一、7（面积）}
填空：520平方米＝（\hspace{1.8em}）平方分米。\par
\begin{BookAnswer}
填\(52000\)。面积单位进率为100：\[1\text{平方米}=100\text{平方分米}\]\[\begin{aligned}520\text{平方米}&=(520\times100)\text{平方分米}\\&=52000\text{平方分米}\end{aligned}\]
\end{BookAnswer}

% p081_s08
\BookQuestion{一、7（钱数）}
填空：587元3角5分＝（\hspace{1.8em}）元。\par
\begin{BookAnswer}
填\(587.35\)。先将角和分化成元：\[3\text{角}=0.3\text{元}\]\[5\text{分}=0.05\text{元}\]再合并钱数：\[\begin{aligned}587\text{元}3\text{角}5\text{分}&=(587+0.3+0.05)\text{元}\\&=587.35\text{元}\end{aligned}\]
\end{BookAnswer}

% p081_s09
\BookQuestion{一、7（质量一）}
填空：6260克＝（\hspace{1.8em}）千克（\hspace{1.8em}）克。\par
\begin{BookAnswer}
依次填\(6\)、\(260\)。先分出整千克对应的克数：\[6260\text{克}=6000\text{克}+260\text{克}\]\[6000\text{克}=6\text{千克}\]所以\[6260\text{克}=6\text{千克}260\text{克}\]
\end{BookAnswer}

% p081_s10
\BookQuestion{一、7（长度）}
填空：329.4厘米＝（\hspace{1.8em}）米。\par
\begin{BookAnswer}
填\(3.294\)。厘米换米除以\(100\)：\[\begin{aligned}329.4\text{厘米}&=(329.4\div100)\text{米}\\&=3.294\text{米}\end{aligned}\]
\end{BookAnswer}

% p081_s11
\BookQuestion{一、7（质量二）}
填空：275.09吨＝（\hspace{1.8em}）吨（\hspace{1.8em}）千克。\par
\begin{BookAnswer}
依次填\(275\)、\(90\)。先分成整数吨与小数吨：\[275.09\text{吨}=275\text{吨}+0.09\text{吨}\]小数吨换成千克：\[\begin{aligned}0.09\text{吨}&=(0.09\times1000)\text{千克}\\&=90\text{千克}\end{aligned}\]所以\[275.09\text{吨}=275\text{吨}90\text{千克}\]
\end{BookAnswer}

% p081_s12
\BookQuestion{一、8}
把（\hspace{1.8em}）的小数点向左移动两位是86.27；把（\hspace{1.8em}）扩大到原来的1000倍是7460。\par
\begin{BookAnswer}
依次填\(8627\)、\(7.46\)。左移两位后为\(86.27\)，原数应为\(86.27\times100=8627\)；扩大\(1000\)倍后为\(7460\)，原数\(7460\div1000=7.46\)。
\end{BookAnswer}

% p081_s13
\BookQuestion{一、9}
一个两位小数，保留一位小数是236.5，这个数最小是（\hspace{1.8em}），最大是（\hspace{1.8em}）。\par
\begin{BookAnswer}
依次填\(236.45\)、\(236.54\)。原数是两位小数，保留一位要看百分位。\par 最小情况：十分位原来是\(4\)，百分位至少为\(5\)才能进成\(236.5\)，最小是\(236.45\)。\par 最大情况：十分位原来是\(5\)，百分位至多为\(4\)才不会进成\(236.6\)，最大是\(236.54\)。\(236.44\)会舍成\(236.4\)，\allowbreak{}\(236.55\)会进成\(236.6\)，都不符合。
\end{BookAnswer}

% p081_s14
\BookQuestion{一、10}
学校图书馆原有6.26万册图书，上周借出1.3万册且现在未还回，本周又新进2.25万册图书，现在图书馆有（\hspace{1.8em}）万册图书。\par
\begin{BookAnswer}
填\(7.21\)。按题中当前在馆本数计算：\(6.26-1.3+2.25=4.96+2.25=7.21\)(万册)。借出未还先减去，新进图书再加上。
\end{BookAnswer}

\Needspace{8\baselineskip}

\pdfbookmark[2]{原书第82页}{orig-82}

\subsection*{原书第82页}

% p082_s01
\BookQuestion{一、11}
11. 一个数，百位和百分位上都是\(6\)，其余数位上都是\(0\)，这个数是（\hspace{1.8em}）。\par
\begin{BookAnswer}
填\(600.06\)。百位的\(6\)表示\(600\)，百分位的\(6\)表示\(0.06\)，其余数位为\(0\)，合为\(600+0.06=600.06\)。
\end{BookAnswer}

% p082_s02
\BookQuestion{二、1}
二、选择（将正确答案的序号填在括号里）。（每题2分，共12分）\par 1. 用简便方法计算\(602\times32\)时，下面方法正确的是（\hspace{1.8em}）。\par
\noindent A. \(32\times600+32\)\par
\noindent B. \(32\times600+32\times2\)\par
\noindent C. \(32+600+32\times2\)\par
\begin{BookAnswer}
选B。利用乘法交换律、分配律，把\(602\)拆成\(600+2\)：\[\begin{aligned}602\times32&=32\times(600+2)\\&=32\times600+32\times2\end{aligned}\]A只加了\(32\)，少加一个\(32\)；C把应有的\(32\times600\)错写为\(32+600\)。
\end{BookAnswer}

% p082_s03
\BookQuestion{二、2}
2. 下面与\(287\)最接近的是（\hspace{1.8em}）。\par
\noindent A. 286.999\par
\noindent B. 287.01\par
\noindent C. 287.1\par
\begin{BookAnswer}
选A。分别求与\(287\)的差：\par \(287-286.999=0.001\)；\par \(287.01-287=0.01\)；\par \(287.1-287=0.1\)。\par 差\(0.001\)最小，所以\(286.999\)最近。
\end{BookAnswer}

% p082_s04
\BookQuestion{二、3}
3. 一个小数的小数点先向右移动两位，又向左移动三位后是\(327.06\)，原来的小数是（\hspace{1.8em}）。\par
\noindent A. 32.706\par
\noindent B. 327.06\par
\noindent C. 3270.6\par
\begin{BookAnswer}
选C。先右移两位再左移三位，合起来是向左移一位，即原数除以\(10\)。原数应为\(327.06\times10=3270.6\)。
\end{BookAnswer}

% p082_s05
\BookQuestion{二、4}
4. \(24\times[623\div(12+34)-15]\)的运算顺序是（\hspace{1.8em}）。\par
\noindent A. 乘→除→加→减\par
\noindent B. 加→除→减→乘\par
\noindent C. 除→加→减→乘\par
\begin{BookAnswer}
选B。先算小括号\(12+34\)的加法，再算\(623\)除以这个和，然后减\(15\)，最后与\(24\)相乘，顺序为“加、除、减、乘”。
\end{BookAnswer}

% p082_s06
\BookQuestion{二、5}
5. \(24\)个\(32\)相加的和除以\(12\)的\(22\)倍，商是多少？列式为（\hspace{1.8em}）。\par
\noindent A. \(24\times32\div12\times22\)\par
\noindent B. \(24\times32\div(12\times22)\)\par
\noindent C. \((24+32)\div(12\times22)\)\par
\begin{BookAnswer}
选B。\(24\)个\(32\)相加是\(32\times24\)，利用乘法交换律也可写为\(24\times32\)；\allowbreak{}\(12\)的\(22\)倍是\(12\times22\)。求前者除以后者，除数必须把整个乘积括起来：\[(32\times24)\div(12\times22)\]这与B等值。A没有括号，按从左到右计算，就变成先除以\(12\)再乘\(22\)；C把“相加的和”误写成\(24+32\)。
\end{BookAnswer}

% p082_s07
\BookQuestion{二、6}
6. 在学校运动会的跳高比赛中，王兵跳了\(1.23\)米，比张宇高出了\(4\)厘米，张宇跳了（\hspace{1.8em}）米。\par
\noindent A. 1.27\par
\noindent B. 1.53\par
\noindent C. 1.19\par
\begin{BookAnswer}
选C。先统一单位：\[4\text{厘米}=0.04\text{米}\]王兵比张宇高\(0.04\)米，故张宇\(1.23-0.04=1.19\)米。
\end{BookAnswer}

% p082_s08
\BookQuestion{三、1（1）—（4）}
三、计算。（共26分）\par 1. 直接写得数。（8分）\par \(37\times2\times5=\quad 25\times101=\quad20.03+6.7=\quad46.25+8.6=\)\par
\begin{BookAnswer}
\(37\times2\times5=37\times10=370\)；\allowbreak{}\(25\times101=25\times100+25=2525\)；\allowbreak{}\(20.03+6.7=26.73\)；\allowbreak{}\(46.25+8.6=54.85\)。小数加法把同数位相加，必要时补\(0\)对齐。
\end{BookAnswer}

% p082_s09
\BookQuestion{三、1（5）—（8）}
1. 直接写得数（续）。\par \(278-32-18=\quad50.8+62.7=\quad182-23.6=\quad55.4+9.46=\)\par
\begin{BookAnswer}
\(278-32-18=278-(32+18)=228\)；\allowbreak{}\(50.8+62.7=113.5\)；\allowbreak{}\(182-23.6=158.4\)；\allowbreak{}\(55.4+9.46=64.86\)。小数点对齐后计算。
\end{BookAnswer}

% p082_s10
\BookQuestion{三、2（1）}
2. 计算下列各题，能简算的要简算。（18分）\par \(233.2+56.4+26.8+83.6\)\par
\begin{BookAnswer}
\textbf{凑整求和}\[\begin{aligned}&233.2+56.4+26.8+83.6\\&=(233.2+26.8)+(56.4+83.6)\\&=260+140\\&=400\end{aligned}\]
\end{BookAnswer}

% p082_s11
\BookQuestion{三、2（2）}
\(200-63.46-24.54\)\par
\begin{BookAnswer}
\textbf{合并减数}\[\begin{aligned}&200-63.46-24.54\\&=200-(63.46+24.54)\\&=200-88\\&=112\end{aligned}\]
\end{BookAnswer}

% p082_s12
\BookQuestion{三、2（3）}
\(87\times102\)\par
\begin{BookAnswer}
把\(102\)拆成\(100+2\)，用乘法分配律计算。\[\begin{aligned}&87\times102\\={}&87\times(100+2)\\={}&87\times100+87\times2\\={}&8700+174\\={}&8874\end{aligned}\]
\end{BookAnswer}

\Needspace{8\baselineskip}

\pdfbookmark[2]{原书第83页}{orig-83}

\subsection*{原书第83页}

% p083_s01
\BookQuestion{三、2（4）}
2. 计算下列各题，能简算的要简算（续）。\par \(540-36\times43\div9\)\par
\begin{BookAnswer}
\textbf{先算整除的一组}\par 乘积除以\(9\)，可先把其中能整除的\(36\)除以\(9\)，再乘\(43\)。\[\begin{aligned}&540-36\times43\div9\\&=540-(36\div9)\times43\\&=540-4\times43\\&=540-172\\&=368\end{aligned}\]
\end{BookAnswer}

% p083_s02
\BookQuestion{三、2（5）}
\(25\times32\times250\)\par
\begin{BookAnswer}
\textbf{拆\(32\)凑整}\[\begin{aligned}&25\times32\times250\\&=(25\times4)\times(8\times250)\\&=100\times2000\\&=200000\end{aligned}\]
\end{BookAnswer}

% p083_s03
\BookQuestion{三、2（6）}
\(246\div[96\div(854-838)]\)\par
\begin{BookAnswer}
\textbf{先小括号再中括号}\[\begin{aligned}&246\div[96\div(854-838)]\\&=246\div[96\div16]\\&=246\div6\\&=41\end{aligned}\]
\end{BookAnswer}

% p083_s04
\BookQuestion{四、1}
四、解决问题。（共31分）\par 1. 一辆小汽车和一辆货车分别从烟台和北京同时出发，相向而行，经过\(3\)小时，两车还相距\(193\)千米。已知小汽车的速度为\(95\)千米/时，货车的速度为\(75\)千米/时，烟台和北京相距多少千米？（先画图表示数量关系，再解答）（5分）\par
\begin{BookAnswer}
两车已行的路程加上还相隔的\(193\)千米，就是两地的距离。\par
\textbf{方法一　先求两车的速度和}
\[\begin{aligned}&(95+75)\times3+193\\={}&170\times3+193\\={}&510+193\\={}&703\text{(千米)}\end{aligned}\]
答：烟台和北京相距\(703\)千米。\par
\textbf{方法二　先分别求两车的路程}
\[\begin{aligned}&95\times3+75\times3+193\\={}&285+225+193\\={}&510+193\\={}&703\text{(千米)}\end{aligned}\]
答：烟台和北京相距\(703\)千米。
\BookFigure{figures/p083_s04-01.tex}
\end{BookAnswer}

% p083_s05
\BookQuestion{四、2}
2. 周末王叔叔去超市采购生活物资。他买了一袋\(5\)千克装的大米，花了\(79.5\)元，又买了一桶\(5\)升装的食用油，买大米比买食用油少花了\(23.9\)元。王叔叔买大米和食用油一共花了多少元？（6分）\par
\begin{BookAnswer}
大米花了\(79.5\)元，食用油比大米多花\(23.9\)元。用加法交换律和结合律，先把两个\(79.5\)相加。
\[\begin{aligned}&79.5+23.9+79.5\\={}&(79.5+79.5)+23.9\\={}&159+23.9\\={}&182.9\text{(元)}\end{aligned}\]
答：买大米和食用油一共花了\(182.9\)元。
\end{BookAnswer}

% p083_s06
\BookQuestion{四、3}
3. 为了迎接植树节，学校组织了“绿色校园”植树劳动实践活动。三年级有\(158\)名同学参加，四年级参加的人数比三年级的\(2\)倍少\(68\)人，五年级有\(362\)名同学参加，五年级参加的人数比四年级多多少人？（6分）\par
\begin{BookAnswer}
先求三年级人数的\(2\)倍，再减去\(68\)人得到四年级人数。用五年级人数减去四年级人数，求出相差的人数。
\[\begin{aligned}&362-(158\times2-68)\\={}&362-(316-68)\\={}&362-248\\={}&114\text{(人)}\end{aligned}\]
答：五年级参加的人数比四年级多\(114\)人。
\end{BookAnswer}

\Needspace{8\baselineskip}

\pdfbookmark[2]{原书第84页}{orig-84}

\subsection*{原书第84页}

% p084_s01
\BookQuestion{四、4（1）}
4. 四季水果店即将推出“周末特惠”活动，李老板提前备货，购入葡萄和橙子各\(35\)箱。葡萄每箱\(55\)元，橙子每箱\(45\)元。\par （1）购入葡萄和橙子一共花了多少钱？（4分）\par
\begin{BookAnswer}
两种水果各35箱，可把一箱葡萄与一箱橙子配成一组，先求每组价钱。也可分别按单价乘箱数列出总价，再用乘法分配律把55与45凑成100。\par \[\begin{aligned}&55\times35+45\times35\\={}&(55+45)\times35\\={}&100\times35\\={}&3500\text{(元)}\end{aligned}\]答：一共花3500元。
\end{BookAnswer}

% p084_s02
\BookQuestion{四、4（2）}
（2）葡萄比橙子多花了多少钱？（4分）\par
\begin{BookAnswer}
两种水果箱数相同，每箱价差乘35箱就是总价差。也可先列两种总价之差，再用乘法分配律把55减45凑成10。\par \[\begin{aligned}&55\times35-45\times35\\={}&(55-45)\times35\\={}&10\times35\\={}&350\text{(元)}\end{aligned}\]答：葡萄比橙子多花350元。
\end{BookAnswer}

% p084_s03
\BookQuestion{四、5}
5. 王老师为学校新一年级的同学采购桌椅，不小心把明细表弄脏了，你能帮他算出每把椅子多少元吗？（6分）\par
\par\smallskip\Needspace{4\baselineskip}\begingroup\small\setlength{\tabcolsep}{3pt}
\noindent\begin{tabularx}{\linewidth}{|*{4}{>{\centering\arraybackslash}X|}}\hline
物品 & 单价 & 数量 & 总价\\\hline
课桌 & 125 & 6 & 1014\\\hline
椅子 & 被污渍遮挡 & 12 & 同上（合计）\\\hline
\end{tabularx}\par
\noindent 1014是两项合计；原表总价单元格跨两行。\par
\endgroup\smallskip
\begin{BookAnswer}
题表的\(1014\)元是桌、椅两项总价。利用“椅子单价\(=\)椅子总价\(\div\)椅子数量”，必须先从两项总价中扣掉课桌费用。\textbf{方法一：总价扣除课桌费用，再除以椅子数量}\[\begin{aligned}&(1014-125\times6)\div12\\&=(1014-750)\div12\\&=264\div12\\&=22\text{(元)}\end{aligned}\]答：每把椅子\(22\)元。\textbf{方法二：按一张课桌配两把椅子分组}\par 椅子\(12\)把、课桌\(6\)张，可分成\(6\)组，每组一桌两椅。\[\begin{aligned}&(1014\div6-125)\div2\\&=(169-125)\div2\\&=44\div2\\&=22\text{(元)}\end{aligned}\]答：每把椅子\(22\)元。
\end{BookAnswer}

% p084_s04
\BookQuestion{五、思考题}
五、动脑筋想一想。（共6分）\par 王军和张兵去书店买书，他们看中了一套两本装的《趣味数学》。王军说：“我的钱不够，还差9.3元。”张兵说：“我的钱不够，还差14.7元。”\par 后来两人把钱合在一起，正好买下了这套书。假设这套书中两本书的价钱相同，你知道平均每本书多少钱吗？\par
\begin{BookAnswer}
两人的钱合起来恰好买一套，所以王军差的钱就是张兵已有的钱，张兵差的钱就是王军已有的钱。一套书的价钱等于两人各自所差的钱之和；每套有两本且单价相同。\textbf{先求一套价钱，再求一本价钱}\[\begin{aligned}&(9.3+14.7)\div2\\&=24\div2\\&=12\text{(元)}\end{aligned}\]答：每本书\(12\)元。
\end{BookAnswer}

\Needspace{12\baselineskip}

\section{总复习（二）}

\Needspace{8\baselineskip}

\pdfbookmark[2]{原书第85页}{orig-85}

\subsection*{原书第85页}

% p085_s01
\BookQuestion{一、1}
总复习（二）（图形与几何）。时间：90分钟；满分：100分（知识：98分，卷面：2分）。\par 一、填空。（每空1分，共23分）\par 1. 一条直线的平行线有（\hspace{1.8em}）条，过直线外一点可以画（\hspace{1.8em}）条直线与这条直线垂直。\par
\begin{BookAnswer}
依次填无数、\(1\)。一条直线可有无数条不同距离的平行线；在同一平面内，过直线外同一点只能作一条垂线。
\end{BookAnswer}

% p085_s02
\BookQuestion{一、2}
2. 长方形的邻边相互（\hspace{1.8em}），对边相互（\hspace{1.8em}）。\par
\begin{BookAnswer}
依次填垂直、平行。长方形的四个内角都是直角，相邻两边互相垂直；两组对边分别平行。
\end{BookAnswer}

% p085_s03
\BookQuestion{一、3}
3. 以平行四边形的一条边为底，能作出（\hspace{1.8em}）条高，这些高的长度（\hspace{1.8em}）。\par
\begin{BookAnswer}
依次填无数、相等。以一条边为底，从对边任意一点向这条底作垂线段，都得到这个底上的高，所以有无数条。它们都是同一组平行线之间的垂线段，两条平行线间距离处处相等，因此同一底上的高都相等；不能把另一组底所对应的高混在一起比较。
\end{BookAnswer}

% p085_s04
\BookQuestion{一、4}
4. 右图中，有（\hspace{1.8em}）个三角形，有（\hspace{1.8em}）个平行四边形，有（\hspace{1.8em}）个梯形。\par
\BookFigure{figures/p085_s04-01.tex}
\begin{BookAnswer}
依次填\(5,4,6\)。把上方从左到右四点记为\(A,B,C,D\)，下方从左到右三点记为\(E,F,G\)。\par 三角形共\(5\)个：\[\begin{array}{ccc}ABE&BEF&BCF\\CFG&CDG&\end{array}\]平行四边形共\(4\)个：\[\begin{array}{cc}ABFE&ACGE\\BCFE&BCGF\end{array}\]恰有一组对边平行的梯形共\(6\)个：\[\begin{array}{cc}ACFE&ADGE\\BCGE&BDGE\\BDGF&CDGF\end{array}\]
\end{BookAnswer}

% p085_s05
\BookQuestion{一、5}
5. 已知一个三角形的两条边长分别是\(5\,\mathrm{cm}\)和\(9\,\mathrm{cm}\)，第三条边最长是（\hspace{1.8em}）\(\mathrm{cm}\)，最短是（\hspace{1.8em}）\(\mathrm{cm}\)。（取整厘米数）\par
\begin{BookAnswer}
依次填\(13\)、\(5\)。三角形任意两边之和都要大于第三边。因此第三边要小于\(5+9=14\)，又要满足\(5+x>9\)，即大于\(9-5=4\)。所以\(4<x<14\)，取整厘米时最长\(13\)厘米、最短\(5\)厘米。取\(4\)或\(14\)会使两条短边的和等于长边，只能排成线段，不能围成三角形。
\end{BookAnswer}

% p085_s06
\BookQuestion{一、6}
6. 如图，在直角三角形\(ABC\)中（单位：厘米），\allowbreak{}\(BC\)边上的高是（\hspace{1.8em}）厘米，\allowbreak{}\(AB\)边上的高是（\hspace{1.8em}）厘米。\par
\BookFigure{figures/p085_s06-01.tex}
\begin{BookAnswer}
依次填\(2.4\)、\(3\)。以\(BC\)为底，从顶点\(A\)到\(BC\)的垂线段是\(AD\)，所以高是\(2.4\)厘米。以\(AB\)为底，\allowbreak{}\(AC\perp AB\)，所以高是\(AC\)，长\(3\)厘米。
\end{BookAnswer}

% p085_s07
\BookQuestion{一、7（1）}
7. 下面哪些是平移现象，哪些是旋转现象？\par （1）张叔叔在公路上开车，方向盘的运动是（\hspace{1.8em}）现象。\par
\begin{BookAnswer}
填旋转。方向盘绕固定的转轴转动，各点围绕同一个中心改变方向。
\end{BookAnswer}

% p085_s08
\BookQuestion{一、7（2）}
（2）升国旗时，国旗的上升运动是（\hspace{1.8em}）现象。\par
\begin{BookAnswer}
填平移。国旗整体沿旗杆向上移动，形状、大小和朝向不变。
\end{BookAnswer}

% p085_s09
\BookQuestion{一、7（3）}
（3）妈妈用拖布拖地，是（\hspace{1.8em}）现象。\par
\begin{BookAnswer}
填平移。按题中沿直线推拉的运动理解，拖布整体沿同一方向移动。
\end{BookAnswer}

% p085_s10
\BookQuestion{一、7（4）}
（4）小明骑自行车，车轮的运动是（\hspace{1.8em}）现象。\par
\begin{BookAnswer}
填旋转。车轮绕车轴转动。这里判断的是“车轮转动”，不是自行车整体前进。
\end{BookAnswer}

% p085_s11
\BookQuestion{一、8}
8. 正方形有（\hspace{1.8em}）条对称轴，等边三角形有（\hspace{1.8em}）条对称轴，等腰梯形有（\hspace{1.8em}）条对称轴。\par
\begin{BookAnswer}
依次填\(4,3,1\)。正方形的4条对称轴是两条对角线所在的直线和两组对边中点连线所在的直线；等边三角形有3条；通常所说的等腰梯形有经过两底中点的1条。
\end{BookAnswer}

% p085_s12
\BookQuestion{一、9（钟面）}
9. 钟面显示的是（\hspace{1.8em}）时，将分针顺时针旋转〔原扫描中个数不清〕个直角后是（\hspace{1.8em}）时（\hspace{1.8em}）分。\par
\BookFigure{figures/p085_s12-01.tex}
\begin{BookAnswer}
钟面可以读作\(3\)时。原扫描中可辨认“直角”，但旋转几个直角的数字不清，后两个空暂不能唯一确定。\par 仅在该处为“1个直角”，即\(90^\circ\)的条件下：钟面每一大格对应\(30^\circ\)，分针转过的大格数是\[\begin{aligned}&90\div30\\={}&3\text{(大格)}\end{aligned}\]每大格表示5分钟，经过的时间按“每格分钟数\(\times\)格数”计算：\[\begin{aligned}&5\times3\\={}&15\text{(分钟)}\end{aligned}\]所以在这个条件下，后两个空填3、15，即3时15分。原题的旋转个数核清前，不能把这个条件解当作无条件答案。
\end{BookAnswer}

\Needspace{8\baselineskip}

\pdfbookmark[2]{原书第86页}{orig-86}

\subsection*{原书第86页}

% p086_s01
\BookQuestion{二、1}
二、选择（将正确答案的序号填在括号里）。（每题3分，共18分）\par 1. 将图形顺时针旋转\(90^\circ\)后可以得到下面图形（\hspace{1.8em}）。\par
\noindent A. 三格横排，最右格上方再一格\par
\noindent B. 三格竖排，最下格右侧再一格\par
\noindent C. 三格竖排，最上格左侧再一格\par
\noindent D. 无法确定\par
\BookFigure{figures/p086_s01-01.tex}
\BookFigure{figures/p086_s01-02.tex}
\BookFigure{figures/p086_s01-03.tex}
\BookFigure{figures/p086_s01-04.tex}
\begin{BookAnswer}
选C。顺时针转\(90^\circ\)后，原来的三格横排变成竖排，左下伸出的那一格转到上端左侧，与C相同。旋转只改变方向，不能反射翻面。
\end{BookAnswer}

% p086_s02
\BookQuestion{二、2}
2. 在同一平面内的\(a,b,c\)三条直线，直线\(a\)与直线\(b\)都和直线\(c\)相交成直角，则直线\(a\)与直线\(b\)（\hspace{1.8em}）。\par
\noindent A. 互相平行\par
\noindent B. 互相垂直\par
\noindent C. 不平行也不垂直\par
\noindent D. 无法确定\par
\begin{BookAnswer}
选A。同一平面内，两条直线都垂直于同一条直线，这两条直线互相平行，即\(a\perp c\)、\(b\perp c\)推出\(a\parallel b\)。
\end{BookAnswer}

% p086_s03
\BookQuestion{二、3}
3. 等腰三角形中，有一个内角是\(40^\circ\)，另外两个内角是（\hspace{1.8em}）。\par
\noindent A. \(40^\circ\)和\(100^\circ\)\par
\noindent B. \(70^\circ\)和\(50^\circ\)\par
\noindent C. \(40^\circ\)和\(100^\circ\)或者\(70^\circ\)和\(70^\circ\)\par
\noindent D. 无法确定\par
\begin{BookAnswer}
选C。按给定内角\(40^\circ\)：若它是顶角，另两角均为\((180-40)\div2=70^\circ\)；若它是底角，另一底角也是\(40^\circ\)，顶角\(180-40\times2=100^\circ\)。两种情况都要考虑。
\end{BookAnswer}

% p086_s04
\BookQuestion{二、4}
4. 把一个长方形框架拉成一个平行四边形，平行四边形的周长与原长方形的周长（\hspace{1.8em}）。\par
\noindent A. 大\par
\noindent B. 小\par
\noindent C. 一样大\par
\noindent D. 无法比较大小\par
\begin{BookAnswer}
选C。铁丝框拉成平行四边形，四条边的长度没有变，只是角度和高改变，所以周长不变，与原长方形一样大。
\end{BookAnswer}

% p086_s05
\BookQuestion{二、5}
5. 下列图形中，不是轴对称图形的是（\hspace{1.8em}）。\par
\noindent A. 向右箭头\par
\noindent B. 一般平行四边形\par
\noindent C. 十字形轮廓\par
\noindent D. 圆形禁止标志\par
\BookFigure{figures/p086_s05-01.tex}
\begin{BookAnswer}
选B。A沿水平中线对折能重合，C沿水平或竖直中线对折都能重合。D沿斜杠所在的直线，或沿过圆心且垂直于斜杠的直线对折，都能重合。B的角不是直角，相邻两边长度不等，这个平行四边形没有对称轴。
\end{BookAnswer}

% p086_s06
\BookQuestion{二、6}
6. 以下是某手机地图APP中的四个工具标识图，哪个不是轴对称图形？（\hspace{1.8em}）\par
\noindent A. 家人地图\par
\noindent B. 驾车巡航\par
\noindent C. 新能源导航\par
\noindent D. 卫星求救\par
\BookFigure{figures/p086_s06-01.tex}
\begin{BookAnswer}
选C。A、B的外轮廓和内部图案分别沿各自的竖直中线对折，都能重合。C的车身轮廓左右对称，车内的闪电图案沿竖直中线对折不能重合。D沿穿过机身和天线中心的斜线对折，太阳能板和信号弧线都能重合。
\end{BookAnswer}

% p086_s07
\BookQuestion{三、1}
三、判断（对的打“√”，错的打“×”）。（共6分）\par 1. 过直线外一点只能画一条与已知直线方向平行的直线。（\hspace{1.8em}）\par
\begin{BookAnswer}
判\(\checkmark\)。在同一平面内，过直线外一点，有且只有一条直线与已知直线平行。这是平行线的唯一性。
\end{BookAnswer}

% p086_s08
\BookQuestion{三、2}
2. 不相交的两条直线互相平行。（\hspace{1.8em}）\par
\begin{BookAnswer}
判\(\times\)。必须补充“在同一平面内”。空间里两条直线可能既不相交，也不平行，例如不同方向且不在同一平面的两条棱所在直线。
\end{BookAnswer}

% p086_s09
\BookQuestion{三、3}
3. 一个等腰三角形，底角是[疑读\(26^\circ\)，原扫描末位不清]，它一定是钝角三角形。（\hspace{1.8em}）\par
\begin{BookAnswer}
判断为\(\checkmark\)。若底角为\(26^\circ\)，顶角\(180^\circ-26^\circ\times2=128^\circ\)，是钝角三角形；若原字确为\(25^\circ\)，顶角则为\(130^\circ\)，结论同样为钝角。原扫描末位近似6，不能把130度作为已核清的唯一顶角；判断√不受这两种读法影响。
\end{BookAnswer}

% p086_s10
\BookQuestion{三、4}
4. 旋转后图形的大小、形状都改变了。（\hspace{1.8em}）\par
\begin{BookAnswer}
判\(\times\)。旋转前后形状、大小都不变，改变的是位置和朝向。
\end{BookAnswer}

% p086_s11
\BookQuestion{三、5}
5. 相交后的四个角都是直角的两条直线一定互相垂直。（\hspace{1.8em}）\par
\begin{BookAnswer}
判\(\checkmark\)。相交形成的四个角都是直角，正是两条直线互相垂直的定义，所以可判垂直。
\end{BookAnswer}

% p086_s12
\BookQuestion{三、6}
6. 两个等底等高的梯形一定能拼成一个平行四边形。（\hspace{1.8em}）\par
\begin{BookAnswer}
判\(\times\)。等底等高不保证两个梯形的腰长、倾斜方式都相同，未必能拼成两组对边分别平行的四边形。能够保证拼成平行四边形的是两个完全相同的梯形，沿合适的腰拼接；不能把“等底等高”当成“形状、大小完全相同”。
\end{BookAnswer}

\Needspace{8\baselineskip}

\pdfbookmark[2]{原书第87页}{orig-87}

\subsection*{原书第87页}

% p087_s01
\BookQuestion{四、1（1）}
四、操作。（共27分）\par 1.（1）画出左图的另一半，使它成为一个轴对称图形。（4分）\par
\BookFigure{figures/p087_s01-01.tex}
\begin{BookAnswer}
以虚线为对称轴，逐个顶点数到轴的水平格数，再在轴另一侧同一水平线上取相同距离的点，按原轮廓顺次连接。对称轴上的点保持原位，对应连线垂直于对称轴；不能只凭形状相像。左页红线按这一点对点规则补出。
\BookFigure{figures/p087_s01-02.tex}
\end{BookAnswer}

% p087_s02
\BookQuestion{四、1（2）}
（2）将右图绕点\(O\)顺时针旋转\(90^\circ\)，再向左平移\(3\)格。（5分）\par
\BookFigure{figures/p087_s02-01.tex}
\begin{BookAnswer}
先绕\(O\)顺时针旋转\(90^\circ\)，再向左平移\(3\)格。先确认旋转的三个要素：中心\(O\)、顺时针方向、角度\(90^\circ\)。\par 旋转时\(O\)不动，其他顶点到\(O\)的距离不变；从\(O\)出发的每条对应线段顺时针转成直角。再按“向左、\(3\)格”将旋转所得图形的三个顶点各移相同格数，依次相连。第二步针对旋转后的图形，不能直接平移原图。
\BookFigure{figures/p087_s02-02.tex}
\end{BookAnswer}

% p087_s03
\BookQuestion{四、2（1）}
2.（1）将图形向上平移\(3\)格后，再向右平移\(5\)格。（4分）\par
\BookFigure{figures/p087_s03-01.tex}
\begin{BookAnswer}
标①的红图是原三角形先向上\(3\)格、再向右\(5\)格后的图。每个顶点都按同一平移过程移动，两条直角边的方向、长度不变。
\BookFigure{figures/p087_s03-02.tex}
\end{BookAnswer}

% p087_s04
\BookQuestion{四、2（2）}
（2）将图形绕点\(O\)顺时针旋转\(90^\circ\)。（4分）\par
\BookFigure{figures/p087_s04-01.tex}
\begin{BookAnswer}
标②的红图是原三角形绕\(O\)顺时针旋转\(90^\circ\)后的图。\(O\)保持不动，原来朝左的边转到朝上，另一个顶点应随同一方向转动，不能把整个图平移到别处。
\BookFigure{figures/p087_s04-02.tex}
\end{BookAnswer}

% p087_s05
\BookQuestion{四、2（3）}
（3）以已知直线为对称轴作对称图形。（4分）\par
\BookFigure{figures/p087_s05-01.tex}
\begin{BookAnswer}
标③的红图是关于已知竖直虚线的轴对称图。原图左侧各顶点映到轴右侧，距离相同；轴上的\(O\)保持原位。三次操作分别针对原图，不是把①接着变成②、③。
\BookFigure{figures/p087_s05-02.tex}
\end{BookAnswer}

% p087_s06
\BookQuestion{四、3}
3. 分别画出下面图形底边上的高。（6分）\par
\BookFigure{figures/p087_s06-01.tex}
\begin{BookAnswer}
三角形：从上方顶点向标“底”的下边作垂线段。平行四边形：从上边顶点向下底作垂线段。梯形：从下底端点向标“底”的上边作垂线段。\par 使用三角尺时，一条直角边靠住底所在直线，沿底移动到需要经过的顶点，再沿另一条直角边画出高。每条高都须与对应底垂直，垂足位于底所在直线上，直角标记贴住垂足；斜腰不能直接当作高。
\BookFigure{figures/p087_s06-02.tex}
\end{BookAnswer}

% p087_s07
\BookQuestion{五、1（1）}
五、解决问题。（共24分）\par 1.（1）请你画一条从王庄到加油站最近的路。（3分）\par
\BookFigure{figures/p087_s07-01.tex}
\begin{BookAnswer}
王庄与加油站作为两个位置点，直接连接这两点，得到最短路线；依据是两点之间线段最短。左页“到加油站”的红线连接的是这两个位置。
\BookFigure{figures/p087_s07-02.tex}
\end{BookAnswer}

% p087_s08
\BookQuestion{五、1（2）}
（2）请你画一条从王庄到公路最近的路。（3分）\par
\BookFigure{figures/p087_s08-01.tex}
\begin{BookAnswer}
从王庄向靠近王庄的公路边线作垂线段，垂足就在该边线上，直角标记须贴住交点。这条“到公路”的红线最短，依据是点到直线的所有连线中垂线段最短。两道“最短”所对应的对象不同。
\BookFigure{figures/p087_s08-02.tex}
\end{BookAnswer}

\Needspace{8\baselineskip}

\pdfbookmark[2]{原书第88页}{orig-88}

\subsection*{原书第88页}

% p088_s01
\BookQuestion{五、2}
五、解决问题（续）。\par 2. 王师傅要制作一个周长是\(68\)厘米的等腰梯形框架，其中上底是\(8\)厘米，下底是\(16\)厘米，则这个等腰梯形的一条腰长是多少厘米？（6分）\par
\begin{BookAnswer}
等腰梯形两腰相等。\textbf{方法一：先扣除两底，再平均分给两腰}\[\begin{aligned}&[68-(8+16)]\div2\\={}&[68-24]\div2\\={}&44\div2\\={}&22\text{(厘米)}\end{aligned}\]答：一条腰长22厘米。\par\textbf{方法二：从一半周长扣除两底和的一半}\[\begin{aligned}&68\div2-(8+16)\div2\\={}&34-24\div2\\={}&34-12\\={}&22\text{(厘米)}\end{aligned}\]答：一条腰长22厘米。
\end{BookAnswer}

% p088_s02
\BookQuestion{五、3}
3. 一个平行四边形的一条边长\(24\)厘米，比它的邻边短\(2\)厘米，这个平行四边形的周长是多少厘米？（6分）\par
\begin{BookAnswer}
平行四边形两组对边分别相等，四条边是两条\(24\)厘米和两条\(24+2=26\)厘米；求周长要把四边相加，不能把邻边误当成也长\(24\)厘米。\textbf{方法一：先求邻边，再算周长}\[\begin{aligned}&[24+(24+2)]\times2\\&=[24+26]\times2\\&=50\times2\\&=100\text{(厘米)}\end{aligned}\]答：平行四边形周长\(100\)厘米。\textbf{方法二：先按四条\(24\)厘米，再补两条边各多的\(2\)厘米}\[\begin{aligned}&24\times4+2\times2\\&=96+4\\&=100\text{(厘米)}\end{aligned}\]答：周长\(100\)厘米。
\end{BookAnswer}

% p088_s03
\BookQuestion{五、4}
4. 小明用一根铁丝围成边长是\(6\)厘米的正方形，如果要围成一个等边三角形，那么它的边长是多少厘米？（6分）\par
\begin{BookAnswer}
\textbf{总长不变，列式后先算能整除的一组}\par 按每边6厘米乘4条边求原总长，再除以3。乘积除以3可先用6除以3，再乘4。\[\begin{aligned}&6\times4\div3\\={}&6\div3\times4\\={}&2\times4\\={}&8\text{(厘米)}\end{aligned}\]答：等边三角形的边长是8厘米。
\end{BookAnswer}

\Needspace{12\baselineskip}

\section{总复习（三）}

\Needspace{8\baselineskip}

\pdfbookmark[2]{原书第89页}{orig-89}

\subsection*{原书第89页}

% p089_s01
\BookQuestion{一、1}
总复习（三）（统计与概率）。时间：40分钟；满分：100分（知识：98分，卷面：2分）。\par 一、填空。（每空3分，共24分）\par 1.（\hspace{1.8em}）能较好地反映一组数据的整体水平。\par
\begin{BookAnswer}
填平均数。平均数把一组数据的总量平均分配，能反映这组数据的整体水平，但不一定等于其中某一个实际数据。
\end{BookAnswer}

% p089_s02
\BookQuestion{一、2}
2. 小玉一分钟跳绳\(160\)个，小涵一分钟跳绳\(150\)个，他俩的平均跳绳个数是（\hspace{1.8em}）个，小涵再多跳（\hspace{1.8em}）个，他俩的跳绳数就一样多了。\par
\begin{BookAnswer}
依次填155、10。简便法以150为基准，把多出的10个平均补给2人：\[\begin{aligned}&150+(160-150)\div2\\={}&150+10\div2\\={}&150+5\\={}&155\text{(个)}\end{aligned}\]另法，两人总次数除以2：\[\begin{aligned}&(160+150)\div2\\={}&310\div2\\={}&155\text{(个)}\end{aligned}\]小涵补到小玉的160个，须多跳\(160-150=10\)个。
\end{BookAnswer}

% p089_s03
\BookQuestion{一、3}
3. 四个数\(12,8,7,13\)的平均数是（\hspace{1.8em}）。\par
\begin{BookAnswer}
填10。用加法结合律凑两个20：\[\begin{aligned}&(12+8+7+13)\div4\\={}&[(12+8)+(7+13)]\div4\\={}&[20+20]\div4\\={}&40\div4\\={}&10\end{aligned}\]
\end{BookAnswer}

% p089_s04
\BookQuestion{一、4}
4. 四年级\(3\)个班向希望小学捐赠图书，平均每班捐\(20\)本，其中，一班捐了\(23\)本，二班捐了\(15\)本，三班捐了（\hspace{1.8em}）本。\par
\begin{BookAnswer}
填22。\textbf{方法一：移多补少}二班比平均数少5本，一班比平均数多3本，三班应比平均数多2本。利用加减运算的性质，先算5减3。\[\begin{aligned}&20+(20-15)-(23-20)\\={}&20+5-3\\={}&20+(5-3)\\={}&20+2\\={}&22\text{(本)}\end{aligned}\]\textbf{方法二：总数扣除已知两班}\[\begin{aligned}&20\times3-(23+15)\\={}&60-38\\={}&22\text{(本)}\end{aligned}\]
\end{BookAnswer}

% p089_s05
\BookQuestion{一、5}
5. 三个连续自然数的平均数是\(12\)，其中最大的数是（\hspace{1.8em}）。\par
\begin{BookAnswer}
填\(13\)。连续三个自然数可以把最大数的\(1\)移给最小数，三数恰好相同，且总和不变，所以平均数等于中间数\(12\)。原来三数为\(11,12,13\)，最大的是\(13\)。
\end{BookAnswer}

% p089_s06
\BookQuestion{一、6}
6. 在仰卧起坐达标测试中，\allowbreak{}\(A\)组\(3\)名同学共做\(96\)个，\allowbreak{}\(B\)组\(4\)名同学共做\(156\)个。平均每人做（\hspace{1.8em}）个。\par
\begin{BookAnswer}
填36。总次数除以总人数：\[\begin{aligned}&(96+156)\div(3+4)\\={}&252\div7\\={}&36\text{(个)}\end{aligned}\]两组人数不同，不能直接平均两组的平均数。
\end{BookAnswer}

% p089_s07
\BookQuestion{一、7}
7. 绘画比赛有\(7\)个评委，分别给张梦打分\(98\)分、\(95\)分、\(96\)分、\(92\)分、\(95\)分、\(96\)分、\(99\)分，按照比赛规则，要去掉一个最高分和一个最低分，其他分数的平均数就是最后得分。张梦的最后得分是（\hspace{1.8em}）分。\par
\begin{BookAnswer}
填96。去掉一个最高分99和一个最低分92，剩5个评分。\textbf{方法一：以96分为基准移多补少}98多2分，两个95各少1分，两次96不变。连续减去两个数，等于减去它们的和：\[\begin{aligned}&96+(2-1-1)\div5\\={}&96+[2-(1+1)]\div5\\={}&96+[2-2]\div5\\={}&96+0\div5\\={}&96+0\\={}&96\text{(分)}\end{aligned}\]\textbf{方法二：分组求总分再平均}用加法交换律、结合律分组：\[\begin{aligned}&(98+95+96+95+96)\div5\\={}&[98+(96+96)+(95+95)]\div5\\={}&[98+192+190]\div5\\={}&[290+190]\div5\\={}&480\div5\\={}&96\text{(分)}\end{aligned}\]答：张梦最后得分是96分。
\end{BookAnswer}

% p089_s08
\BookQuestion{二、1}
二、选择（将正确答案的序号填在括号里）。（每题2分，共10分）\par 1. 一个人工湖测得最深处是\(2.7\)米，最浅处是\(1.1\)米，这个人工湖的平均水深可能是（\hspace{1.8em}）。\par
\noindent A. 2.7米\par
\noindent B. 1.1米\par
\noindent C. 1.9米\par
\noindent D. 0.8米\par
\begin{BookAnswer}
选C。平均水深必须在最浅\(1.1\)米和最深\(2.7\)米之间，且这里最深最浅不同，平均不能等于任一端点。\(1.9\)米符合，\allowbreak{}\(0.8\)米小于最浅处。
\end{BookAnswer}

% p089_s09
\BookQuestion{二、2}
2. 四年级一班学生的平均体重是\(35\)千克，小萱在四年级一班，她的体重（\hspace{1.8em}）是\(32\)千克。\par
\noindent A. 可能\par
\noindent B. 不可能\par
\noindent C. 一定\par
\noindent D. 无法确定\par
\begin{BookAnswer}
选A。平均体重\(35\)千克，并不要求每个人都是\(35\)千克；小萱\(32\)千克完全可能。例如\(32,35,38\)的平均数就是\(35\)。
\end{BookAnswer}

\Needspace{8\baselineskip}

\pdfbookmark[2]{原书第90页}{orig-90}

\subsection*{原书第90页}

% p090_s01
\BookQuestion{二、3}
3. 气象观测站在某天的\(3\)时、\(8\)时、\(15\)时、\(21\)时测得的温度分别是\(8^\circ\mathrm C\)、\(15^\circ\mathrm C\)、\(26^\circ\mathrm C\)、\(19^\circ\mathrm C\)。计算这天的平均气温正确的算式为（\hspace{1.8em}）。\par
\noindent A. \((8+15+26+19)\div4\)\par
\noindent B. \((8+15+26+19)\div(3+8+15+21)\)\par
\noindent C. \((8+15+26+19)\div24\)\par
\noindent D. \((3+8+15+21)\div4\)\par
\begin{BookAnswer}
选A。统计的是\(4\)个观测时刻的温度值\(8,15,26,19\)，所以用温度总和除以观测次数：\[(8+15+26+19)\div4\]时刻\(3,8,15,21\)表示何时观测，不是气温数据；\allowbreak{}\(24\)也不是这组数据的个数，不能作为除数。
\end{BookAnswer}

% p090_s02
\BookQuestion{二、4}
4. 一次跳绳比赛，\allowbreak{}\(5\)名同学的成绩从低到高依次是\(73\)下、\(81\)下、\(a\)下、\(87\)下、\(95\)下，他们的平均成绩可能是（\hspace{1.8em}）。\par
\noindent A. 72下\par
\noindent B. 85下\par
\noindent C. 95下\par
\noindent D. 96下\par
\begin{BookAnswer}
按四舍五入保留整数的口径，选B。由从低到高排列得\(81\leq a\leq87\)。先求已知四人的成绩和：
\[\begin{aligned}&73+81+87+95\\={}&(73+87)+(81+95)\\={}&160+176\\={}&336\end{aligned}\]
当\(a=81\)时，平均数最小：
\[\begin{aligned}&(336+81)\div5\\={}&417\div5\\={}&83.4\text{(下)}\end{aligned}\]
当\(a=87\)时，平均数最大：
\[\begin{aligned}&(336+87)\div5\\={}&423\div5\\={}&84.6\text{(下)}\end{aligned}\]
若允许四舍五入保留整数，\allowbreak{}\(84.6\)下约为\(85\)下。题干未说明取整，按精确平均数计算，四个选项均不符合。
\end{BookAnswer}

% p090_s03
\BookQuestion{二、5}
5. \(5\)个数的平均数是\(50\)，若把其中一个数改为\(80\)，则平均数变为\(60\)，这个数原来是（\hspace{1.8em}）。\par
\noindent A. 60\par
\noindent B. 50\par
\noindent C. 40\par
\noindent D. 30\par
\begin{BookAnswer}
选D。改前和改后都是\(5\)个数。平均数从\(50\)变为\(60\)，总和增加了\(5\)个\(10\)。只有一个数改动，从\(80\)减去总和增加的数，就得到原来的数。
\[\begin{aligned}&80-(60-50)\times5\\={}&80-10\times5\\={}&80-50\\={}&30\end{aligned}\]
\end{BookAnswer}

% p090_s04
\BookQuestion{三、1（1）—（4）}
三、计算。（共23分）\par 1. 直接写得数。（8分）\par \(0.4+0.5=\quad1-0.7=\quad4.3+3.7=\quad1.2-0.8=\)\par
\begin{BookAnswer}
\(0.4+0.5=0.9\)；\allowbreak{}\(1-0.7=0.3\)；\allowbreak{}\(4.3+3.7=8\)；\allowbreak{}\(1.2-0.8=0.4\)。把十分之一的个数相加减，满\(10\)个十分之一就进为\(1\)。
\end{BookAnswer}

% p090_s05
\BookQuestion{三、1（5）—（8）}
1. 直接写得数（续）。\par \(0.4+1.2=\quad7.9-1.9=\quad7.3+1.8=\quad13.5-3.5=\)\par
\begin{BookAnswer}
\(0.4+1.2=1.6\)；\allowbreak{}\(7.9-1.9=6\)；\allowbreak{}\(7.3+1.8=9.1\)；\allowbreak{}\(13.5-3.5=10\)。同数位相加减，小数点位置不变。
\end{BookAnswer}

% p090_s06
\BookQuestion{三、2（1）}
2. 列竖式计算。（带☆的要验算）（15分）\par \(15.05+2.63=\textcolor{answerred}{17.68}\)\par
\begin{BookAnswer}
\[\begin{array}{r@{\;}r@{.}l}&15&05\\+&2&63\\\hline&17&68\end{array}\]\(15.05+2.63=17.68\)。百分位\(5+3=8\)，十分位\(0+6=6\)，个位\(5+2=7\)，十位保留\(1\)。
\end{BookAnswer}

% p090_s07
\BookQuestion{三、2（2）}
\(9.06-2.5=\textcolor{answerred}{6.56}\)\par
\begin{BookAnswer}
\[\begin{array}{r@{\;}r@{.}l}&9&06\\-&2&50\\\hline&6&56\end{array}\]\(9.06-2.5=6.56\)。先把\(2.5\)写成等值的\(2.50\)对齐小数点；十分位不够减，向个位借\(1\)。
\end{BookAnswer}

% p090_s08
\BookQuestion{三、2（3）}
\(5.7+12.34=\textcolor{answerred}{18.04}\)\par
\begin{BookAnswer}
\[\begin{array}{r@{\;}r@{.}l}&5&70\\+&12&34\\\hline&18&04\end{array}\]\(5.7+12.34=18.04\)。\(5.7\)补写成\(5.70\)，十分位\(7+3=10\)向个位进\(1\)，所以结果十分位要保留\(0\)。
\end{BookAnswer}

% p090_s09
\BookQuestion{三、2（4）及验算}
☆\(12.8-1.24=\textcolor{answerred}{11.56}\)（要验算）\par
\begin{BookAnswer}
\[\begin{array}{r@{\;}r@{.}l}&12&80\\-&1&24\\\hline&11&56\end{array}\]\(12.8-1.24=11.56\)。百分位不够减，从十分位借\(1\)，再逐位相减。验算用差加减数：\[\begin{array}{r@{\;}r@{.}l}&11&56\\+&1&24\\\hline&12&80\end{array}\]验算得到原被减数\(12.80=12.8\)，计算正确。
\end{BookAnswer}

\Needspace{8\baselineskip}

\pdfbookmark[2]{原书第91页}{orig-91}

\subsection*{原书第91页}

% p091_s01
\BookQuestion{四、1}
四、解决问题。（共41分）\par 1. 一次数学竞赛中\(4\)位同学的成绩分别是\(96\)分、\(88\)分、[第三数扫描缺损，疑读92]分、[第四数扫描缺损，疑读18/88]分，他们的平均成绩是多少分？（7分）\par
\begin{BookAnswer}
第3、4项字形有缺损，不能认定原题已核清。以下只按96、88、92、88计算。\textbf{方法一：分组求总分再平均}先按加法结合律分成两组，才能同一步求两个和：\[\begin{aligned}&(96+88+92+88)\div4\\={}&[(96+88)+(92+88)]\div4\\={}&[184+180]\div4\\={}&364\div4\\={}&91\text{(分)}\end{aligned}\]答：按这组数据，平均成绩是91分。\par\textbf{方法二：以88分为基准移多补少}\[\begin{aligned}&88+[(96-88)+(92-88)]\div4\\={}&88+[8+4]\div4\\={}&88+12\div4\\={}&88+3\\={}&91\text{(分)}\end{aligned}\]答：按这组数据，平均成绩是91分。\par 若缺损两项实际为82、18，则：\[\begin{aligned}&(96+88+82+18)\div4\\={}&[(96+88)+(82+18)]\div4\\={}&[184+100]\div4\\={}&284\div4\\={}&71\text{(分)}\end{aligned}\]答：仅按该备选读法，平均成绩是71分；仍须核清原题。
\end{BookAnswer}

% p091_s02
\BookQuestion{四、2}
2. 学校合唱队男生有\(8\)人，平均身高是\(132\)厘米，女生有\(12\)人，平均身高是\(142\)厘米，这个合唱队的学生平均身高是多少厘米？（7分）\par
\begin{BookAnswer}
\textbf{方法一：身高总和除以总人数}\par 男生与女生人数不同，须先分别求身高总和，再除以对应的总人数。不能把两个小组的平均身高直接相加除以2，因为这样把8名男生与12名女生当成了同样人数的两组。\[\begin{aligned} &(132\times8+142\times12)\div(8+12)\\ ={}&(1056+1704)\div20\\ ={}&2760\div20\\ ={}&138\text{(厘米)}\end{aligned}\]答：合唱队学生平均身高是138厘米。\par \textbf{方法二：以132厘米为基准}\par 把每人先按132厘米计算；12名女生每人多10厘米，再平均补给20人。\[\begin{aligned} &132+(142-132)\times12\div(8+12)\\ ={}&132+10\times12\div20\\ ={}&132+120\div20\\ ={}&132+6\\ ={}&138\text{(厘米)}\end{aligned}\]答：合唱队学生平均身高是138厘米。\par
\end{BookAnswer}

% p091_s03
\BookQuestion{四、3}
3. 学校开展图书活动，前\(2\)天平均每天借书\(160\)本，后\(6\)天一共借阅\(464\)本，学校平均每天借书多少本？（8分）\par
\begin{BookAnswer}
\textbf{方法一：总本数除以总天数}\par 前2天借书总量是每天160本乘2天，再加后6天的464本。\[\begin{aligned} &(160\times2+464)\div(2+6)\\ ={}&(320+464)\div8\\ ={}&784\div8\\ ={}&98\text{(本)}\end{aligned}\]答：平均每天借书98本。\par \textbf{基准复核：与每天100本比较}\par 主法已求得总量784本。8天若平均每天100本，共应借800本，少的16本平均到8天，每天少2本。\[\begin{aligned} &100-[100\times8-(160\times2+464)]\div8\\ ={}&100-[800-(320+464)]\div8\\ ={}&100-[800-784]\div8\\ ={}&100-16\div8\\ ={}&100-2\\ ={}&98\text{(本)}\end{aligned}\]答：平均每天借书98本。\par
\end{BookAnswer}

% p091_s04
\BookQuestion{四、4}
4. 甲、乙、丙三人的平均肺活量是[扫描重影，原答案按2500，疑读3000]毫升，甲的肺活量是\(2700\)毫升，乙的肺活量是\(2200\)毫升，丙的肺活量是多少毫升？（8分）\par
\begin{BookAnswer}
平均数有重影，以下仅按2500毫升作条件计算。\textbf{方法一：先合并甲、乙的量}用减法性质合并两个减数：\[\begin{aligned}&2500\times3-(2700+2200)\\={}&7500-4900\\={}&2600\text{(毫升)}\end{aligned}\]答：若平均数是2500毫升，丙的肺活量是2600毫升。\par\textbf{方法二：移多补少}乙少300毫升，甲多200毫升，丙应在平均数上补足100毫升：\[\begin{aligned}&2500+[(2500-2200)-(2700-2500)]\\={}&2500+[300-200]\\={}&2500+100\\={}&2600\text{(毫升)}\end{aligned}\]答：若平均数是2500毫升，丙的肺活量是2600毫升。\par 若原平均数实际为3000毫升，则：\[\begin{aligned}&3000\times3-(2700+2200)\\={}&9000-4900\\={}&4100\text{(毫升)}\end{aligned}\]答：若平均数是3000毫升，丙的肺活量是4100毫升。须核清原题后确定唯一答案。
\end{BookAnswer}

\Needspace{8\baselineskip}

\pdfbookmark[2]{原书第92页}{orig-92}

\subsection*{原书第92页}

% p092_s01
\BookQuestion{四、5（1）}
5. 下面是实验小学四年级二班男生和女生校服订制情况统计表。\par （1）将上面的两个统计表合并成一个统计表。（6分）\par
\par\smallskip\Needspace{4\baselineskip}\begingroup\small\setlength{\tabcolsep}{3pt}
\noindent\textbf{四年级二班女生校服订制情况统计表}\par\smallskip
\noindent\begin{tabularx}{\linewidth}{|*{5}{>{\centering\arraybackslash}X|}}\hline
型号 & 小号 & 中号 & 大号 & 特大号\\\hline
人数 & 2 & 12 & 4 & 3\\\hline
\end{tabularx}\par
\endgroup\smallskip
\par\smallskip\Needspace{4\baselineskip}\begingroup\small\setlength{\tabcolsep}{3pt}
\noindent\textbf{四年级二班男生校服订制情况统计表}\par\smallskip
\noindent\begin{tabularx}{\linewidth}{|*{5}{>{\centering\arraybackslash}X|}}\hline
型号 & 小号 & 中号 & 大号 & 特大号\\\hline
人数 & 3 & 12 & 7 & 2\\\hline
\end{tabularx}\par
\endgroup\smallskip
\par\smallskip\Needspace{4\baselineskip}\begingroup\small\setlength{\tabcolsep}{3pt}
\noindent\textbf{四年级二班男生和女生校服订制情况统计表}\par\smallskip
\noindent\begin{tabularx}{\linewidth}{|*{6}{>{\centering\arraybackslash}X|}}\hline
型号 & 小号 & 中号 & 大号 & 特大号 & 合计/人\\\hline
女生人数 &  &  &  &  & \\\hline
男生人数 &  &  &  &  & \\\hline
总计/人 &  &  &  &  & \\\hline
\end{tabularx}\par
\endgroup\smallskip
\begin{BookAnswer}
把两张表按相同型号合并，同一型号的女生人数与男生人数相加。
\[\begin{array}{c|rrrr|r}
\text{型号}&\text{小号}&\text{中号}&\text{大号}&\text{特大号}&\text{合计/人}\\\hline
\text{女生人数}&2&12&4&3&21\\
\text{男生人数}&3&12&7&2&24\\
\text{总计/人}&5&24&11&5&45
\end{array}\]
\textbf{女生合计}
\[\begin{aligned}&2+12+4+3\\={}&(2+3)+(12+4)\\={}&5+16\\={}&21\text{(人)}\end{aligned}\]
\textbf{男生合计}
\[\begin{aligned}&3+12+7+2\\={}&(3+7)+(12+2)\\={}&10+14\\={}&24\text{(人)}\end{aligned}\]
小号\(2+3=5\)人，中号\(12+12=24\)人，大号\(4+7=11\)人，特大号\(3+2=5\)人。全班共有\(21+24=45\)人。
\end{BookAnswer}

% p092_s02
\BookQuestion{四、5（2）}
（2）女生校服订制（\hspace{1.8em}）的人数最少，男生校服订制（\hspace{1.8em}）的人数最多。（4分）\par
\begin{BookAnswer}
女生各型号订制人数是\(2,12,4,3\)，\allowbreak{}\(2\)最小，对应小号。男生各型号订制人数是\(3,12,7,2\)，\allowbreak{}\(12\)最大，对应中号。\par
答：女生订制小号的人数最少，男生订制中号的人数最多。
\end{BookAnswer}
